\honest{Absolute Authority}{Eliezer Yudkowsky}{2008}

Someone tells you that science ``doesn't really know anything,'' because scientists changed their minds about gravity, so ``who's to say that tomorrow you won't change your minds about evolution?'' \nb{The question contains an argument from history: past theories were replaced, so present ones may be too. Philosophers of science call it the pessimistic induction. The post answers the demand for certainty and does not return to gravity. The Asimov essay that Yudkowsky had quoted the day before, at the end of ``The Fallacy of Gray,'' answers it: good theories are refined, ``not so much wrong as incomplete.''}

For such people, whom I call ``the unenlightened ones,'' a source is trusted wholly or not at all, and a proponent who admits a flaw has confessed. ``One obvious source'' of this habit is religion, where doubting scripture is a sin. ``I suspect'' school is another, since there the teacher's statements must be believed. I add, in passing, that people who want an ``Authoritative probability'' before they will use Bayesian probability are like this too, since someone ``might \ldots\ argue with your estimate of the prior probability.'' \nb{That worry is a real problem in statistics, called the problem of the priors, and it divides Bayesians themselves.}

Such people, I say, toss the word ``certainty'' around, while scientists predict the Moon's next rising and still call it a probabilistic guess. I have them say ``I'm perfect.'' \nb{No one is quoted saying it.}

In a debate I might try a few lines. Science's strength is changing its mind. ``Not absolutely certain'' in science means a tiny chance, not a specific doubt. And if you would change your mind about something for enough evidence, ``it can't have a probability exactly equal to one.'' \nb{That is correct, and statisticians call it Cromwell's rule.}

In private, the first lesson is that you can make moral and factual distinctions without certainty. Some say that without certainty ``the moral relativists win.'' But the error in the relativists' chain is the step from gray to all one shade of gray, and certainty is sufficient for choosing, not necessary. That is my main point.

The learner also needs to know that probability means calibration. Ask a hundred people for statements they are ``absolutely certain'' of, and not all hundred will be right. \nb{That is true. In one study quoted in Yudkowsky's own chapter on global risks, answers given odds of a million to one or more were right 90\% of the time.} ``If anything,'' fanatical beliefs are less often right than obvious ones, so emotional commitment ``doesn't even behave monotonically.'' \nb{The post gives no data for this. In the study just mentioned, accuracy rose with stated odds at the top of the scale, though stated odds are not quite the commitment the post means.}

Probability 1.0 would be infinite certainty. So, ``maybe,'' scientists should say ``We are not infinitely certain.''
